\documentclass[11pt,a4paper]{article}

\usepackage[utf8]{inputenc}
\usepackage[T1]{fontenc}
\usepackage{lmodern}
\usepackage{graphicx}
\usepackage{amsmath,amssymb}
\usepackage{microtype}
\usepackage{booktabs}
\usepackage{siunitx}
\usepackage[numbers,sort&compress]{natbib}
\usepackage[hidelinks]{hyperref}
\usepackage{geometry}
\geometry{margin=1in}
\usepackage{authblk}
\usepackage{titlesec}
\titleformat{\section}{\normalfont\Large\bfseries}{\thesection}{1em}{}
\titleformat{\subsection}{\normalfont\large\bfseries}{\thesubsection}{1em}{}

\title{Extracellular Potassium Relays Axonal Activity to Oligodendrocytes via Kir4.1 to Sustain Axonal Metabolism}
\author{Anonymous}
\date{}

\begin{document}
\maketitle

\begin{abstract}
\noindent
Myelinating oligodendrocytes (OLs) accelerate action potential propagation and supply metabolic fuel to the axons they ensheathe, yet the acute signal that informs mature OLs about ongoing axonal firing remains unclear. Combining two-photon Ca$^{2+}$ imaging with compound action potential (CAP) recordings in acutely isolated mouse optic nerves, we show that high-frequency axonal stimulation triggers robust, frequency-scaled Ca$^{2+}$ surges in mature OLs. The responses require action potentials (abolished by tetrodotoxin, $93 \pm 9\%$ reduction) and extracellular Ca$^{2+}$ ($96 \pm 14\%$ reduction), but depend only marginally on ionotropic glutamate or purinergic receptors and are essentially insensitive to voltage-gated Ca$^{2+}$ channels. Block of oligodendroglial inwardly rectifying K$^+$ channels with 100~\textmu M Ba$^{2+}$ reduces the 50~Hz-evoked OL Ca$^{2+}$ rise by $84 \pm 10\%$ and the $\Delta[\text{K}^+]_\text{bath}$-evoked rise by $88 \pm 9\%$, while the reverse-mode Na$^+$/Ca$^{2+}$ exchanger blockers KB-R7943 and SEA0400 suppress OL Ca$^{2+}$ signals by $\approx$40\%. OL-specific deletion of Kir4.1 (Kir4.1\textsuperscript{fl/fl};MOGiCre) slows CAP recovery after high-frequency activity, reduces MCT1 by $50 \pm 13\%$ and GLUT1 by $44 \pm 8\%$ in central myelin, and lowers basal and activity-induced axonal glucose consumption by $\approx$40\%. Young Kir4.1 cKO axons display reduced resting lactate, blunted stimulus-evoked lactate surges, and impaired ATP recovery after chemical ischaemia and 50~Hz stimulation; axonal pathology emerges at 7--8 months. We propose that extracellular K$^+$ acts as a fast, energy-efficient axon-to-OL signal that, through Kir4.1, gates reverse-mode NCX Ca$^{2+}$ entry and the metabolite-transporter machinery required to sustain axonal function during demanding firing patterns.
\end{abstract}

\section{Introduction}

Myelinating oligodendrocytes (OLs) in the central nervous system do more than accelerate action potential propagation: they also provide trophic and metabolic support to the axons they ensheathe \citep{Nave2010,Simons2016,Philips2017}. The discovery that myelin-resident monocarboxylate transporter~1 (MCT1) delivers lactate to energy-demanding axons, and that shifting myelin toward glycolysis maintains axonal integrity, redefined CNS myelin as a metabolic compartment \citep{Funfschilling2012,Lee2012}. Oligodendroglial N-methyl-D-aspartate (NMDA) receptors then emerged as triggers for surface translocation of GLUT1 and enhanced glucose import during axonal activity \citep{Saab2016}, and genetically encoded sensors confirmed that axonal ATP levels in intact white matter are dynamically coupled to oligodendroglial metabolism \citep{Trevisiol2017,Chamberlain2021}. Despite these advances, the acute signal by which mature OLs sense the moment-to-moment firing state of their axons has remained elusive.

Every action potential extrudes K$^+$ into the extracellular space \citep{Kofuji2000,Olsen2015}, producing periaxonal K$^+$ transients that are energetically cheap, spatially confined, and well-suited to serve as a fast carrier of the axonal firing code. Inwardly rectifying K$^+$ channels of the Kir4.1 (KCNJ10) family are highly expressed on oligodendrocytes and astrocytes and establish the hyperpolarised macroglial resting potential that both follows and buffers these K$^+$ rises \citep{Neusch2001,Djukic2007,Larson2018,Hawkins2017}. Oligodendrocyte-specific deletion of Kir4.1 produces late-onset axonopathy \citep{Schirmer2018}, and loss-of-function of this channel has been implicated in demyelinating and neurodegenerative disease \citep{Tong2014,Stumpf2019,Battefeld2019}. Whether Kir4.1 actively couples ongoing axonal activity to oligodendroglial signalling and to the metabolic support machinery of myelin has, however, not been directly tested.

Here we combine two-photon Ca$^{2+}$ imaging in mature OLs, OL-specific Kir4.1 knockout, pharmacology, quantitative proteomics of optic nerve, and genetically encoded ATP (ATeam1.03), lactate (Laconic) and glucose (FLII12Pglu) sensors to test the hypothesis that extracellular K$^+$ acts as an acute axon-to-OL messenger. We find that mature OLs respond to high-frequency axonal firing with robust Ca$^{2+}$ surges whose magnitude scales with frequency and is dominated by K$^+$-driven depolarisation through Kir4.1 and subsequent reverse-mode NCX Ca$^{2+}$ entry \citep{Stys1992,Tekkok2005,Micu2018}. Conditional ablation of oligodendroglial Kir4.1 slows activity-dependent white-matter K$^+$ clearance, decreases the abundance of MCT1 and GLUT1 in central myelin, reduces axonal lactate and activity-induced glucose consumption, and ultimately produces late-onset axonopathy. These data identify extracellular K$^+$ and oligodendroglial Kir4.1 as a core axis that couples axonal firing to the metabolic support of white-matter axons.

\section{Related Work}

\subsection{Oligodendrocyte-to-axon metabolic coupling}

Genetic strategies that eliminate myelin-specific MCT1 or shift myelin toward glycolysis revealed that OLs export lactate to fuel active axons \citep{Funfschilling2012,Lee2012}. In vitro, glucose and lactate availability are sufficient to modulate OL differentiation and myelination \citep{Rinholm2011}. Oligodendroglial NMDA receptors then emerged as triggers for activity-dependent surface delivery of GLUT1 and for glycolytic support \citep{Saab2016}. Real-time FRET imaging demonstrated that axonal ATP in white matter depends on oligodendroglial glycolysis \citep{Trevisiol2017}, while SIRT2 transfer from myelin further tunes axonal mitochondrial respiration \citep{Chamberlain2021}. These studies establish a framework in which axonal activity drives OL metabolism, but leave open the identity of the earliest obligatory signal that reaches mature OLs during routine firing.

\subsection{Potassium dynamics and Kir4.1 in glia}

Kir4.1 sets a hyperpolarised resting potential that both tracks and buffers extracellular K$^+$ transients \citep{Kofuji2000,Olsen2015,Hawkins2017}. Genetic removal of Kir4.1 from astrocytes or oligodendrocyte-lineage cells impairs K$^+$ and glutamate clearance, produces early lethality in global knockouts, and causes progressive white-matter defects when restricted to oligodendrocytes \citep{Djukic2007,Neusch2001,Larson2018,Schirmer2018}. Barium blockade of Kir channels slows K$^+$ clearance in optic nerve and alters CAP recovery after high-frequency firing \citep{Bay2012,Larson2018}. Downregulation of Kir4.1 accompanies Huntington's disease and multiple sclerosis pathology \citep{Tong2014,Schirmer2018,Stumpf2019}, placing this channel at the intersection of homeostasis and disease.

\subsection{Ionic mechanisms of OL Ca$^{2+}$ signalling}

Somatic and myelinic Ca$^{2+}$ rises in OLs have been attributed to ionotropic glutamate receptors \citep{Micu2006,Saab2016}, purinergic receptors \citep{Kettenmann2013}, and voltage-gated Ca$^{2+}$ channels. Reverse-mode Na$^+$/Ca$^{2+}$ exchanger (NCX) operation has long been implicated in CNS white matter, contributing to anoxic Ca$^{2+}$ overload and axonal injury \citep{Stys1992,Micu2018}. Axo-myelinic signalling models propose local bidirectional communication at the periaxonal space \citep{Micu2018}. However, the relative contribution of K$^+$/Kir4.1 depolarisation versus neurotransmitter-driven mechanisms to acute OL Ca$^{2+}$ rises during physiological firing has not been quantified in intact white matter, motivating the present study.

\section{Results}

\subsection{High-frequency axonal firing elicits Ca$^{2+}$ surges in mature oligodendrocytes}

\textbf{Context.} Whether mature OLs acutely sense axonal firing has remained unresolved, in part because in situ Ca$^{2+}$ imaging with concurrent electrophysiology in mature OLs is technically demanding.

\textbf{Approach.} We expressed GCaMP6s in mature (CC1\textsuperscript{+}) OLs by treating 6--8-week-old PLP-CreERT;RCL-GCaMP6s mice with tamoxifen and imaged acute optic nerves 4--12 weeks later. Baseline 0.4~Hz firing was punctuated with 30~s trains at 10, 25 or 50~Hz while the compound action potential (CAP) was recorded in parallel with OL somatic Ca$^{2+}$.

\textbf{Finding.} Stimulation elicited robust, frequency-graded Ca$^{2+}$ surges in OLs ($n=108$ cells, $N=8$ nerves), with the Ca$^{2+}$ area under the curve (AUC) rising monotonically from 10 to 25 to 50~Hz (all pairwise $p<0.0001$, one-way ANOVA with Tukey's test). Tetrodotoxin (TTX, 1~\textmu M) abolished the 50~Hz-evoked rise by $93 \pm 9\%$ ($n=82$, $N=6$; $p<0.0001$, paired $t$-test), confirming a strict requirement for propagating action potentials. Removing extracellular Ca$^{2+}$ and chelating residual Ca$^{2+}$ with 200~\textmu M EGTA diminished the response by $96 \pm 14\%$ ($n=44$, $N=3$; $p<0.0001$). The axonal lactate rise measured with Laconic showed a matching frequency scaling (50 vs.\ 25~Hz $p=0.0141$; 50 vs.\ 10~Hz $p<0.0001$; 25 vs.\ 10~Hz $p=0.0141$).

\textbf{Interpretation.} Mature OLs detect axonal firing rate in real time, and the signal is carried by transmembrane Ca$^{2+}$ influx rather than intracellular store release. The parallel frequency dependence of axonal lactate output (Figure~\ref{fig:freq_dependence}) indicates that OL Ca$^{2+}$ and metabolic delivery to axons are co-regulated by the same firing parameter.

\begin{figure}[t]
\centering
\includegraphics[width=0.65\linewidth]{figures/fig4_frequency_dependence.pdf}
\caption{Frequency dependence of the OL Ca$^{2+}$ surge and the axonal lactate rise evoked by 30~s trains at 10, 25 and 50~Hz. Values are normalised to the 50~Hz condition; all pairwise comparisons reached $p<0.05$.}
\label{fig:freq_dependence}
\end{figure}

\subsection{OL Ca$^{2+}$ responses depend on Kir channels and reverse-mode NCX}

\textbf{Context.} Acute Ca$^{2+}$ signalling in OLs has classically been attributed to ionotropic glutamate receptors, purinergic receptors and voltage-gated Ca$^{2+}$ channels, but the quantitative contribution of each source during physiological firing in intact white matter is unknown.

\textbf{Approach.} We systematically dissected candidate Ca$^{2+}$ pathways with selective pharmacology while recording 50~Hz-evoked OL Ca$^{2+}$ surges, and we independently activated OLs by raising extracellular [K$^+$].

\textbf{Finding.} Ionotropic glutamate and purinergic signalling contributed only modestly: 50~\textmu M NBQX left responses unchanged, whereas combined NMDA blockade (100~\textmu M D-AP5 + 100~\textmu M 7-CKA) reduced the 50~Hz-evoked Ca$^{2+}$ rise by approximately 20\%. The broad-spectrum P2X/P2Y antagonists PPADS ($21 \pm 6\%$; $n=46$, $N=3$, $p=0.0018$) and suramin ($22 \pm 6\%$; $n=33$, $N=2$, $p=0.0009$) produced similarly small reductions. Voltage-gated Ca$^{2+}$ channel blockade with cadmium, nickel, nifedipine, benidipine or ruthenium red had at most minor effects (Figure~\ref{fig:drug_effects}).

In striking contrast, 100~\textmu M Ba$^{2+}$, which preferentially blocks Kir channels, reversibly reduced the 50~Hz-evoked OL Ca$^{2+}$ surge by $84 \pm 10\%$ ($n=45$, $N=4$; $p<0.0001$, paired $t$-test). Raising the extracellular K$^+$ to 5, 10 and 30~mM evoked Ca$^{2+}$ rises in OLs that scaled with [K$^+$] (5 vs.\ 10~mM $p=0.0048$; 5 vs.\ 30~mM $p<0.0001$; 10 vs.\ 30~mM $p<0.0001$; one-way ANOVA; $n=35$--57 cells, $N=4$--5 nerves); these K$^+$-evoked Ca$^{2+}$ responses persisted in TTX ($p=0.8144$), confirming they are a direct consequence of OL depolarisation rather than secondary axonal firing. Ba$^{2+}$ also diminished the $\Delta[\text{K}^+]_\text{bath}$-evoked rise by $88 \pm 9\%$ ($n=47$, $N=3$; $p<0.0001$). Together these results identify Kir-mediated depolarisation as the dominant step linking axonal firing to OL Ca$^{2+}$ influx.

Ouabain-evoked Ca$^{2+}$ transients in OLs were suppressed by 25~\textmu M KB-R7943 ($p<0.0001$), confirming a functional reverse-mode NCX in optic nerve OLs. The same concentration of KB-R7943 reduced the 50~Hz-evoked Ca$^{2+}$ surge by $44 \pm 11\%$ ($n=64$, $N=5$; $p=0.0002$) and the K$^+$-evoked Ca$^{2+}$ rise by $47 \pm 8\%$ ($n=52$, $N=3$; $p<0.0001$). A structurally distinct NCX inhibitor, 10~\textmu M SEA0400, likewise reduced the response by $38 \pm 7\%$ ($n=54$, $N=4$; $p<0.0001$), while 50~\textmu M bumetanide had no effect, ruling out a Na$^+$/K$^+$/Cl$^-$ cotransporter NKCC1-mediated volume artefact.

\textbf{Interpretation.} During physiological firing, K$^+$-driven Kir depolarisation, rather than neurotransmitter receptors or voltage-gated Ca$^{2+}$ channels, carries most of the OL Ca$^{2+}$ signal in optic nerve. Reverse-mode NCX accounts for approximately half of that influx, establishing a K$^+$~$\rightarrow$~Kir~$\rightarrow$~depolarisation~$\rightarrow$~NCX chain as the dominant acute coupler of axonal firing to OL Ca$^{2+}$.

\begin{figure}[t]
\centering
\includegraphics[width=\linewidth]{figures/fig1_drug_effects.pdf}
\caption{Pharmacological dissection of the 50~Hz stimulation-evoked OL Ca$^{2+}$ response. Bars show the percentage reduction of the stimulus-evoked Ca$^{2+}$ AUC upon application of each drug (or 30~mM K$^+$ bath challenge where indicated). Values are drawn from paired recordings across 3--6 optic nerves per condition.}
\label{fig:drug_effects}
\end{figure}

\subsection{OL-specific Kir4.1 deletion slows K$^+$ clearance and causes late-onset axonopathy}

\textbf{Context.} Pharmacological Kir blockade is informative but not selective. We therefore asked whether Kir4.1 expressed specifically on OLs is required for activity-dependent K$^+$ clearance and for long-term axonal integrity.

\textbf{Approach.} We generated OL-specific Kir4.1 knockouts (Kir4.1\textsuperscript{fl/fl};MOGiCre, cKO) and compared them with Kir4.1\textsuperscript{fl/fl} littermate controls (ctrl). At 2--3 months, CAP peak latencies and stimulus-response relationships were indistinguishable between genotypes ($p=0.7637$ for peak~1, $p=0.9958$ for peak~2, $p=0.9265$ for peak~3, one-way ANOVA with Holm-\v{S}{\'i}d{\'a}k's test; two-way ANOVA for stimulus-response $F(10,150) = 0.4445$, $p=0.9224$), and electron microscopy of optic nerves from 3-month-old mice revealed no differences in $g$-ratio or axon-calibre distribution ($n=4$--5 mice; $p=0.1584$ for $g$-ratio; $F(25,175)=1.028$, $p=0.4333$ for size distribution).

However, high-frequency challenges exposed a clear defect in activity-dependent K$^+$ handling. After a 1~min 50~Hz train the CAP peak amplitude dropped to a similar extent in both genotypes ($p=0.2486$, one-way ANOVA), but the subsequent recovery of CAP amplitude was profoundly slower in cKO than in ctrl ($p<0.0001$), mirroring the effect of 100~\textmu M Ba$^{2+}$ in wildtype nerves ($p<0.0001$). CAP peak latency followed the same pattern: its stimulus-evoked increase was larger in cKO ($p=0.0232$) and its recovery slower ($p<0.0001$), and equivalent deficits in recovery kinetics were observed at 25~Hz and 10~Hz (CAP amplitude $p=0.0203$ at 25~Hz, $p=0.0242$ at 10~Hz). These data demonstrate that oligodendroglial Kir4.1 shapes white-matter K$^+$ clearance and conduction recovery at physiologically relevant firing rates.

At 7--8 months, electron microscopy of cKO optic nerves revealed a significant increase in ultrastructural features of axonal injury and degeneration ($p=0.0169$, $n=5$ mice, 10 images each), as well as giant axonal swellings that were unique to cKO. GFAP-immunopositive area was also elevated at 7 months ($p=0.0295$), indicating secondary astrogliosis, while IBA1 labelling was unchanged.

\textbf{Interpretation.} Oligodendroglial Kir4.1 is selectively required for activity-dependent K$^+$ clearance in white matter, not for gross myelin formation. Its chronic loss produces progressive axonopathy that emerges months after the physiological K$^+$-handling defect is first detectable, establishing a temporal hierarchy in which physiology fails before structure.

\subsection{Kir4.1 loss from OLs reduces MCT1 and GLUT1 in CNS myelin}

\textbf{Context.} A potassium-handling defect that precedes axonopathy by months cannot cause axonal damage directly; an intermediate molecular consequence must accrue in myelin itself.

\textbf{Approach.} We performed quantitative proteomics of optic nerves from 2.5-month-old cKO and ctrl mice, long before overt axonal damage, and validated candidate hits by immunoblot of biochemically purified central myelin. \textbf{Finding.} The most strongly reduced proteins in cKO ranked by log$_2$ fold change included Kir4.1 itself, along with oligodendroglial metabolite transporters. Immunoblots of biochemically purified central myelin (Figure~\ref{fig:myelin_transporters}) confirmed the proteomic hits: Kir4.1 abundance was reduced by $81 \pm 13\%$ ($n=3$ per genotype, $p=0.0038$), MCT1 by $50 \pm 13\%$ ($p=0.0179$), and GLUT1 by $44 \pm 8\%$ ($p=0.0044$), while GLUT3 was unchanged ($p=0.2470$; moderated $t$-test, $n=4$ mice per genotype).

\textbf{Interpretation.} The selective preservation of GLUT3 argues against non-specific remodelling of the myelin proteome. Instead, oligodendroglial Kir4.1 specifically controls the abundance of the two transporters most directly responsible for activity-induced metabolite delivery from OLs to axons, identifying MCT1 and GLUT1 as plausible molecular intermediates linking K$^+$ handling to axonal health.

\begin{figure}[t]
\centering
\includegraphics[width=0.7\linewidth]{figures/fig2_myelin_transporters.pdf}
\caption{Quantification of metabolite transporters in biochemically purified central myelin from Kir4.1 cKO and littermate controls. Kir4.1, MCT1 and GLUT1 are reduced in cKO myelin, while GLUT3 is unchanged.}
\label{fig:myelin_transporters}
\end{figure}

\subsection{Axonal ATP recovery is modestly impaired in young Kir4.1 cKO mice}

\textbf{Context.} If MCT1 and GLUT1 loss represent the intermediate step from K$^+$ handling to axonopathy, young cKO axons should already show signs of a metabolic imbalance in vivo.

\textbf{Approach.} We expressed the ATP FRET sensor ATeam1.03 in optic nerve axons via intravitreal AAV and monitored axonal ATP during glucose deprivation plus 5~mM NaN$_3$ mitochondrial inhibition (GD+MI), a chemical-ischaemia-like challenge, and during 50~Hz activity. Basal axonal ATP was comparable between genotypes ($n=6$ per group, $p=0.61$, unpaired $t$-test), as were the decline rate upon GD+MI ($p=0.44$) and the initial recovery rate after reperfusion with 10~mM glucose ($p=0.51$). However, the onset of ATP recovery differed ($p=0.03$), and, crucially, 15--20~min after the onset of recovery control nerves had nearly returned to baseline ATP, whereas cKO nerves recovered significantly less ($p=0.0135$). A comparable, milder deficit appeared during 50~Hz stimulation: ATP decline and initial recovery rates were unchanged, but final recovered ATP levels were slightly lower in cKO ($p=0.0498$). \textbf{Interpretation.} These modest metabolic signatures, appearing well before axonal damage becomes visible at 7--8 months, are consistent with a slow erosion of axonal energy homeostasis that begins once oligodendroglial K$^+$ handling is compromised, rather than with acute energy failure.

\subsection{Axonal lactate and glucose metabolism depend on oligodendroglial Kir4.1}

\textbf{Context.} A specific prediction of the MCT1/GLUT1 findings is that axonal lactate delivery and glucose flux should scale with Kir4.1 function, both acutely (via Ba$^{2+}$) and chronically (in cKO).

\textbf{Approach.} We expressed Laconic (lactate) and FLII12Pglu-700u$\delta$6 (glucose) FRET sensors in optic nerve axons, paired imaging with 50~Hz stimulation and with pharmacological or genetic Kir4.1 inactivation, and used cytochalasin~B to read out glycolytic consumption. In wildtype nerves, axonal lactate rose with stimulation frequency (50 vs.\ 25~Hz $p=0.0141$; 50 vs.\ 10~Hz $p<0.0001$; 25 vs.\ 10~Hz $p=0.0141$; one-way ANOVA with Holm-\v{S}{\'i}d{\'a}k's test), and 100~\textmu M Ba$^{2+}$ reduced the 50~Hz-evoked axonal lactate surge by $34 \pm 11\%$ ($n=7$ mice, paired $t$-test, $p=0.0249$), demonstrating that Kir-channel-mediated OL signalling acutely facilitates lactate supply to active axons. Notably, 100~\textmu M Ba$^{2+}$ did not affect stimulation-evoked axonal Ca$^{2+}$ transients (tested at 25 and 50~Hz, $p=0.49$ and $p=0.6859$), ruling out a non-specific axonal action.

In Kir4.1 cKO nerves, both basal axonal lactate ($n=16$ cKO, $n=8$ ctrl, $p=0.0119$) and 50~Hz-evoked lactate surges ($p=0.0174$) were significantly reduced, while the rate of lactate decline during glucose deprivation was unchanged ($p=0.1318$). Basal axonal glucose levels were comparable between genotypes ($n=7$--9, $p=0.2276$), but treatment with iodoacetate to block glycolysis revealed a $36 \pm 13\%$ slower rate of glucose accumulation in cKO ($n=4$--7, $p=0.0235$), and cytochalasin~B-based assessment of glucose consumption showed that basal (0.1~Hz) consumption rates were $37 \pm 17\%$ lower in cKO ($p=0.0384$). During 50~Hz stimulation, control axons accelerated glycolysis to $2.7 \pm 0.6\%$/min while cKO axons failed to increase consumption ($0.1 \pm 0.2\%$/min, $p=0.0002$, Figure~\ref{fig:axonal_glucose}). Because the fold change from 0.1~Hz to 50~Hz was similar between genotypes ($8.5 \pm 2$ vs.\ $8.5 \pm 1$, $p=0.9858$), the glycolytic machinery itself is not defective; rather, overall axonal glucose flux is set at a lower working point in cKO, $\approx$40\% below control at both rest and during activity. Consistently, after 50~Hz trains axonal glucose overshot baseline by $3.2 \pm 0.3\%$ in cKO versus $1.4 \pm 0.4\%$ in ctrl ($p=0.0006$), reflecting reduced consumption relative to supply.

\textbf{Interpretation.} Oligodendroglial Kir4.1 is required for the physiological coupling between axonal firing and axonal lactate delivery, and for the working-point of axonal glucose metabolism. The preserved 50~Hz/0.1~Hz fold change of glycolytic flux indicates that the axonal glycolytic amplifier remains intact; Kir4.1 controls the substrate supply rather than the downstream enzymatic machinery.

\begin{figure}[t]
\centering
\includegraphics[width=\linewidth]{figures/fig3_axonal_glucose.pdf}
\caption{Axonal glucose dynamics in ctrl and Kir4.1 cKO optic nerves. \textbf{Left}, glucose consumption rate measured with cytochalasin~B at basal 0.1~Hz stimulation and during 50~Hz stimulation. \textbf{Right}, post-stimulation overshoot of axonal glucose, which is elevated in cKO.}
\label{fig:axonal_glucose}
\end{figure}

\section{Discussion}

Our experiments establish extracellular K$^+$ as a dominant, fast axon-to-OL signal in adult white matter and identify Kir4.1 channels as both the sensor and the regulator of a coupled metabolic output. The picture that emerges (Figure~\ref{fig:working_model}) is one in which each axonal action potential delivers a small, transient K$^+$ pulse to the periaxonal space, mature OLs sum these pulses through Kir4.1-mediated depolarisation, and the resulting Ca$^{2+}$ entry through reverse-mode NCX serves as the second messenger that links axonal firing to the transporter architecture required to fuel it. We discuss four implications of this axis.

\begin{figure}[t]
\centering
\includegraphics[width=\linewidth]{figures/fig5_working_model.pdf}
\caption{Working model. Periaxonal K$^+$ released by high-frequency firing depolarises oligodendrocytes through Kir4.1. Depolarisation drives reverse-mode NCX Ca$^{2+}$ entry and, on slower timescales, is required for physiological abundance of MCT1 and GLUT1 in central myelin, sustaining activity-induced lactate and glucose supply to axons.}
\label{fig:working_model}
\end{figure}

\paragraph{Potassium as a fast, broadcast axon-to-OL signal.} The glutamatergic and purinergic pathways that have dominated models of axon-OL communication \citep{Micu2006,Saab2016,Kettenmann2013} together accounted for only $\approx$20\% of the 50~Hz-evoked OL Ca$^{2+}$ rise in our recordings, whereas Ba$^{2+}$ removed $\approx$80\% of it and also eliminated the direct response to elevated bath K$^+$. Unlike vesicular release, K$^+$ extrusion occurs at every node of Ranvier during every action potential, does not require specialised release machinery, and propagates rapidly to nearby OL processes. These properties make periaxonal K$^+$ well-suited to report fine-grained axonal firing patterns to mature OLs \citep{Bay2012,Larson2018,Olsen2015}, while glutamate and ATP may fine-tune slower or context-specific components of the response. This parsimonious reassignment of signalling weight is consistent with the preserved myelin structure of 3-month-old cKO nerves despite their marked K$^+$-handling deficit: the Kir4.1-K$^+$ pathway controls the acute physiology of the axon-OL unit, not its gross architecture.

\paragraph{Reverse-mode NCX links OL depolarisation to Ca$^{2+}$ signalling.} Two mechanistically distinct NCX blockers, KB-R7943 and SEA0400, each suppressed approximately half of the 50~Hz- and K$^+$-evoked OL Ca$^{2+}$ rise, and ouabain-evoked Ca$^{2+}$ entry in OLs was equally sensitive to KB-R7943. Voltage-gated Ca$^{2+}$ channel blockers had only minor effects, placing NCX, rather than L/T-type channels, at the heart of OL depolarisation-to-Ca$^{2+}$ transduction \citep{Stys1992,Micu2018}. Because NCX is electrogenic and sensitive to the Na$^+$ gradient, chronic elevation of OL Na$^+$ in Kir4.1 cKO---expected from impaired K$^+$ uptake and altered Na$^+$/K$^+$-ATPase substrate availability---may bias the exchanger further toward reverse mode, potentially amplifying downstream Ca$^{2+}$ transients under pathological conditions. This prediction is accessible to future work using OL-targeted Na$^+$ imaging.

\paragraph{K$^+$-Kir4.1-NCX signalling gates metabolite transporter abundance.} The most consequential finding is that young Kir4.1 cKO mice---long before any overt axonal damage---carry a coordinated $\approx$40--50\% reduction of MCT1 and GLUT1 in central myelin while GLUT3 remains unchanged, and the acute axonal metabolic phenotype matches this ratio: basal and activity-induced glucose consumption are $\approx$37--40\% lower in cKO, stimulus-evoked axonal lactate surges are blunted by $34\%$ in wildtype Ba$^{2+}$-treated nerves and further reduced in cKO, and ATP recovery after chemical ischaemia or 50~Hz activity is impaired by single-digit percentage points that nevertheless reach statistical significance. The quantitative match between transporter abundance and metabolic flux strongly suggests that MCT1/GLUT1 are rate-limiting in vivo and that oligodendroglial K$^+$ sensing sets their working point. Moreover, because the glucose-consumption fold-change upon 50~Hz stimulation was preserved ($\approx$8-fold in both genotypes), the glycolytic amplifier in axons remains intact, and the cKO phenotype is best interpreted as a chronic reduction in the substrate supply rather than a signalling deficit downstream of glycolysis. This reframes oligodendroglial K$^+$ homeostasis, long viewed as a housekeeping function \citep{Olsen2015,Hawkins2017}, as a regulator of the metabolic pipeline that sustains axons.

\paragraph{Implications for neurological disease.} Kir4.1 is downregulated in multiple sclerosis lesions, amyotrophic lateral sclerosis, and Huntington's disease \citep{Tong2014,Schirmer2018,Stumpf2019}, and antibodies against Kir4.1 have been reported in subsets of patients with demyelinating disease. Our results raise the possibility that the earliest pathophysiological signature of oligodendroglial Kir4.1 loss is not an electrical conduction failure but a silent, activity-scaled reduction in the metabolic support of axons that precedes visible axonopathy by months. Therapeutic strategies that rescue Kir4.1 expression or that bypass the downstream MCT1/GLUT1 deficit---for example, by providing alternative substrates such as ketone bodies \citep{Stumpf2019}---merit evaluation in contexts where Kir4.1 is primarily affected. Conversely, the Ba$^{2+}$-insensitive residual K$^+$-evoked Ca$^{2+}$ component, though small, may be carried by other Kir isoforms expressed on OLs \citep{Battefeld2019,Suminaite2019} and could provide partial compensation in chronic disease.

\paragraph{Limitations and outlook.} Our experiments are restricted to acute optic nerve ex vivo, a white-matter tract with relatively homogeneous fibre composition; whether the same Kir4.1-NCX axis operates with identical weighting in grey-matter OLs or in juxtaparanodal compartments remains to be tested. Although we show a striking correlation between MCT1/GLUT1 abundance and axonal glucose flux, we have not directly manipulated transporter levels in a Kir4.1 wildtype background to formally isolate the downstream branch. Finally, the modest ATP phenotype at 3 months suggests that compensation---possibly through astrocytic Kir4.1, astrocyte-derived lactate \citep{Tekkok2005}, or neuronal glycolysis---is initially protective, which raises the question of what tips this balance at 7--8 months to produce overt axonopathy. Long-term in vivo monitoring of K$^+$, Ca$^{2+}$ and ATP in the axon-OL unit, combined with layered genetic rescue of MCT1/GLUT1 in Kir4.1 cKO, will be needed to test these possibilities. Despite these open questions, the present data establish extracellular K$^+$ and oligodendroglial Kir4.1 as a dominant acute coupler of axonal activity to oligodendroglial Ca$^{2+}$ signalling and to the metabolic support machinery of CNS myelin.

\section{Methods}

\paragraph{Animals.} PLP-CreERT;RCL-GCaMP6s mice were treated with tamoxifen for three consecutive days at 6--8 weeks of age and used 4--12 weeks post-injection for OL-targeted GCaMP6s imaging. Immunohistochemistry with anti-GFP and CC1 confirmed that GCaMP6s expression was restricted to mature OLs (quantification on 12 sections from 4 mice). Oligodendrocyte-specific Kir4.1 knockout mice (Kir4.1\textsuperscript{fl/fl};MOGiCre, cKO) were compared to Kir4.1\textsuperscript{fl/fl} (ctrl) littermates. Experiments used 2--3-month-old mice for electrophysiology, proteomics and metabolite imaging, and additionally 7--8-month-old mice for electron-microscopic and immunohistochemical analysis of axonal integrity.

\paragraph{Acute optic nerve preparation and electrophysiology.} Optic nerves were dissected and continuously superfused with ACSF containing 10~mM glucose (and 10~mM lactate where noted for prolonged glucose-deprivation experiments). CAPs were evoked through a suction electrode at a baseline 0.4~Hz (every 2.5~s) and interrogated with 30~s trains at 10, 25 or 50~Hz. CAP peak amplitude, latency, and area were quantified relative to baseline; recovery kinetics were assessed for up to 10~min after each stimulation train.

\paragraph{Two-photon imaging.} GCaMP6s fluorescence in OL somata was imaged in parallel with CAP recordings. Axonal Ca$^{2+}$ dynamics were assessed by expressing GCaMP6s in retinal ganglion cell axons through intravitreal AAV-Cre injection into RCL-GCaMP6s mice, followed by imaging 3--6 weeks later.

\paragraph{Genetically encoded metabolite sensors.} Axonal ATP (ATeam1.03), lactate (Laconic) and glucose (FLII12Pglu-700u$\delta$6) were imaged in optic nerve axons after intravitreal AAV delivery. FRET ratios were normalised to the minima obtained after glucose deprivation plus 5~mM NaN$_3$ (ATP), after 20~min glucose deprivation (lactate), or after 30~min glucose deprivation (glucose). Glycolytic flux was estimated from the slope of cytosolic glucose decline during 20~\textmu M cytochalasin~B application; the transport-independent response was confirmed with 1~mM iodoacetate.

\paragraph{Pharmacology.} Bath-applied compounds included TTX (1~\textmu M), EGTA (200~\textmu M in Ca$^{2+}$-free ACSF), Ba$^{2+}$ (100~\textmu M), NBQX (50~\textmu M), D-AP5 (100~\textmu M), 7-CKA (100~\textmu M), PPADS (50~\textmu M), suramin (50~\textmu M), Cd$^{2+}$ (50--100~\textmu M), Ni$^{2+}$ (200~\textmu M), nifedipine (50~\textmu M), benidipine (10~\textmu M), ruthenium red (10~\textmu M), bumetanide (50~\textmu M), ouabain (500~\textmu M), KB-R7943 (25~\textmu M), SEA0400 (10~\textmu M), NaN$_3$ (5~mM) and cytochalasin~B (20~\textmu M).

\paragraph{Proteomics and myelin biochemistry.} Optic nerves from 2.5-month-old cKO ($n=4$) and ctrl ($n=4$) mice were processed for label-free quantitative proteomics; differentially regulated proteins were ranked by log$_2$ fold change and FDR, with enrichment analysis performed in WebGestalt. Biochemically purified central myelin from 2.5-month-old cKO ($n=3$) and ctrl ($n=3$) was analysed by immunoblot for Kir4.1, MCT1, GLUT1 and GLUT3.

\paragraph{Microscopy and histology.} Electron microscopy was performed on optic nerves from 3-month-old and 7--8-month-old mice; $g$-ratios, axon size distributions, and markers of axonal injury were quantified on randomly sampled images. GFAP and IBA1 immunolabelling was quantified as positive area in optic nerve sections at 2.5 and 7 months.

\paragraph{Statistics.} Data are reported as mean~$\pm$~SEM. Paired and unpaired two-tailed $t$-tests, one-way ANOVA with Tukey's or Holm-\v{S}{\'i}d{\'a}k's post-hoc, and two-way ANOVA were used as indicated. Proteomic comparisons used moderated $t$-tests. $n$ denotes cells or recordings, $N$ denotes optic nerves, unless mice are explicitly specified.

\bibliographystyle{unsrtnat}
\bibliography{references}

\end{document}
